\documentclass{article}
\usepackage{amsmath,amssymb,amsfonts}
\usepackage{algorithm}
\usepackage{algorithmic}
\usepackage{fullpage}
\usepackage{hyperref}
\begin{document}

%%%%%%%%%%%%%%%%%%%%%%%%%%%%%%%%%%%%%%%%%%%%%%%%%%%%%%%%%%%%%%%%%%%%%%%%%%%%%%%
\section*{Algorithm Name}
\textbf{Stochastic Gradient Descent with Warm Restarts (SGDR)}  
The method follows the ordinary SGD update but the stepsize $\eta_t$ is
periodically reset to a large value $\eta_{\max}$ and then decayed
according to a cosine schedule.  The length of the $k$‑th period is denoted
$M_k$ and may increase geometrically.

%%%%%%%%%%%%%%%%%%%%%%%%%%%%%%%%%%%%%%%%%%%%%%%%%%%%%%%%%%%%%%%%%%%%%%%%%%%%%%%
\section*{Mathematical Formulation}
Let $f:\mathbb{R}^d\rightarrow\mathbb{R}$ be the (possibly stochastic) loss
function and let $g_t$ be a stochastic gradient evaluated at iterate $w_t$:
\[
g_t \;=\; \nabla f(w_t;\,\xi_t),\qquad \xi_t\sim\mathcal{D},
\]
where $\mathcal{D}$ is the data distribution.  
The learning‑rate schedule is
\[
\boxed{
\eta_t \;=\; \eta_{\min}
          \;+\; \frac{\eta_{\max}-\eta_{\min}}{2}
                \Bigl(1+\cos\!\Bigl(\pi\frac{t-T_k}{M_k}\Bigr)\Bigr)},
\tag{1}
\]
where $T_k=\sum_{j=0}^{k-1}M_j$ is the starting index of the current
period $k$ (with $M_0$ given).  The schedule satisfies
\[
\eta_{\min}\le \eta_t\le\eta_{\max},\qquad\forall t\ge0 .
\]

The SGDR iteration reads
\[
\boxed{w_{t+1}=w_t-\eta_t\,g_t},\qquad t=0,1,2,\dots .
\tag{2}
\]

Hyper‑parameters:
\begin{itemize}
  \item $\eta_{\max}>0$ : maximal stepsize after a restart,
  \item $\eta_{\min}\ge0$ : minimal stepsize at the end of a period,
  \item $M_0\in\mathbb{N}$ : length of the first period,
  \item $T_{\rm mult}>1$ : multiplicative factor for period lengths,
        i.e. $M_{k+1}=T_{\rm mult}M_k$.
\end{itemize}

%%%%%%%%%%%%%%%%%%%%%%%%%%%%%%%%%%%%%%%%%%%%%%%%%%%%%%%%%%%%%%%%%%%%%%%%%%%%%%%
\section*{Pseudocode}
\begin{algorithm}[H]
\caption{SGDR}
\begin{algorithmic}[1]
\STATE \textbf{Input:} $\eta_{\max},\;\eta_{\min},\;M_0,\;T_{\rm mult},\;w_0$
\STATE Set $k\leftarrow0$, $M_k\leftarrow M_0$, $T_k\leftarrow0$
\FOR{$t=0,1,2,\dots$}
    \STATE Compute stochastic gradient $g_t=\nabla f(w_t;\xi_t)$
    \STATE $\displaystyle\eta_t\leftarrow \eta_{\min}
          +\frac{\eta_{\max}-\eta_{\min}}{2}
            \bigl(1+\cos(\pi\frac{t-T_k}{M_k})\bigr)$ \hfill\textit{// (1)}
    \STATE $w_{t+1}\leftarrow w_t-\eta_t g_t$ \hfill\textit{// (2)}
    \IF{$t=T_k+M_k-1$}
        \STATE $k\leftarrow k+1$
        \STATE $M_k\leftarrow T_{\rm mult}\,M_{k-1}$
        \STATE $T_k\leftarrow t+1$ \COMMENT{restart: reset schedule}
    \ENDIF
\ENDFOR
\end{algorithmic}
\end{algorithm}

%%%%%%%%%%%%%%%%%%%%%%%%%%%%%%%%%%%%%%%%%%%%%%%%%%%%%%%%%%%%%%%%%%%%%%%%%%%%%%%
\section*{Dry Run Example}
Consider the one‑dimensional quadratic
\[
f(w)=(w-3)^2,\qquad \nabla f(w)=2(w-3).
\]
Take deterministic gradients ($g_t=\nabla f(w_t)$) and the following
constants:
\[
\eta_{\max}=0.1,\qquad \eta_{\min}=0,\qquad M_0=3,\qquad T_{\rm mult}=1.
\]
The cosine schedule (1) for the first three steps becomes
\[
\eta_0=0.1,\quad
\eta_1=0.1\!\left(\tfrac{1+\cos(\pi/3)}{2}\right)=0.075,\quad
\eta_2=0.1\!\left(\tfrac{1+\cos(2\pi/3)}{2}\right)=0.025 .
\]

\begin{center}
\begin{tabular}{c|c|c|c}
$t$ & $w_t$ & $g_t=2(w_t-3)$ & $f(w_t)$ \\ \hline
0 & $0$ & $-6$ & $9$ \\
1 & $w_1=0-0.1(-6)=0.6$ & $2(0.6-3)=-4.8$ & $(0.6-3)^2=5.76$ \\
2 & $w_2=0.6-0.075(-4.8)=1.0$ & $2(1.0-3)=-4.0$ & $4$ \\
3 & $w_3=1.0-0.025(-4.0)=1.1$ & $2(1.1-3)=-3.8$ & $3.61$
\end{tabular}
\end{center}
The iterates move toward the minimiser $w^\star=3$ while the stepsize
decays within the period.

%%%%%%%%%%%%%%%%%%%%%%%%%%%%%%%%%%%%%%%%%%%%%%%%%%%%%%%%%%%%%%%%%%%%%%%%%%%%%%%
\section*{Assumptions}
\begin{enumerate}
  \item[{\bf A1}] (\textbf{Smoothness}) $f$ is $L$‑smooth:
  \[
  f(y)\le f(x)+\langle\nabla f(x),y-x\rangle+\frac{L}{2}\|y-x\|^2,
  \qquad\forall x,y\in\mathbb{R}^d .
  \tag{3}
  \]
  \item[{\bf A2}] (\textbf{Unbiased gradient}) For every $t$,
  \[
  \mathbb{E}[g_t\mid w_t]=\nabla f(w_t).
  \tag{4}
  \]
  \item[{\bf A3}] (\textbf{Bounded variance}) There exists $\sigma^2\ge0$ such
  that
  \[
  \mathbb{E}\!\bigl[\|g_t-\nabla f(w_t)\|^2\mid w_t\bigr]\le\sigma^2,
  \qquad\forall t .
  \tag{5}
  \]
  \item[{\bf A4}] (\textbf{Stepsize bound}) The maximal stepsize satisfies
  \[
  \eta_{\max}\le\frac{1}{L}.
  \tag{6}
  \]
\end{enumerate}
All constants $L,\sigma,\eta_{\min},\eta_{\max}$ are known a‑priori.

%%%%%%%%%%%%%%%%%%%%%%%%%%%%%%%%%%%%%%%%%%%%%%%%%%%%%%%%%%%%%%%%%%%%%%%%%%%%%%%
\section*{Main Convergence Theorem}
\begin{theorem}[Convergence of SGDR]\label{thm:sgdr}
Under assumptions {\bf A1}–{\bf A4}, let $\{w_t\}$ be generated by
Algorithm~SGDR for $T$ iterations.  Define
\[
\Delta_f:=f(w_0)-\inf_{w}f(w).
\]
Then the following bound holds:
\[
\frac{1}{T}\sum_{t=0}^{T-1}\mathbb{E}\bigl[\|\nabla f(w_t)\|^2\bigr]
\;\le\;
\underbrace{\frac{2\Delta_f}{\eta_{\min}T}}_{\text{decaying term}}
\;+\;
\underbrace{\frac{L\sigma^2\eta_{\max}}{\eta_{\min}}}_{\text{variance floor}} .
\tag{7}
\]
Consequently, if $\eta_{\max}$ is chosen to decay as
$\eta_{\max}=c/\sqrt{T}$ (e.g. by increasing the period lengths
geometrically), the right–hand side becomes $\mathcal{O}(1/\sqrt{T})$,
which matches the optimal rate for non‑convex stochastic optimisation.
\end{theorem}

%%%%%%%%%%%%%%%%%%%%%%%%%%%%%%%%%%%%%%%%%%%%%%%%%%%%%%%%%%%%%%%%%%%%%%%%%%%%%%%
\section*{Theoretical Properties}
\begin{itemize}
  \item \textbf{Stability.} Because $\eta_t\le\eta_{\max}\le1/L$, the
    coefficient $1-\frac{L}{2}\eta_t\ge\frac12$ in the descent
    inequality, guaranteeing that each iteration does not increase the
    expected objective value.
  \item \textbf{Hyper‑parameter sensitivity.}
    The bound \eqref{7} shows a linear dependence on the ratio
    $\eta_{\max}/\eta_{\min}$.  A too small $\eta_{\min}$ inflates the
    variance floor, while a too large $\eta_{\max}$ may violate
    \eqref{6} and break the descent property.
  \item \textbf{Comparison to plain SGD.}
    For constant stepsize $\eta_t\equiv\eta$, inequality \eqref{7}
    reduces to the classical SGD bound
    $\frac{1}{T}\sum\mathbb{E}\|\nabla f(w_t)\|^2
     \le \frac{2\Delta_f}{\eta T}+L\sigma^2\eta$.
    Warm restarts improve the constant by allowing a larger
    $\eta_{\max}$ early on while still ensuring a diminishing average
    stepsize through the increasing periods.
\end{itemize}

%%%%%%%%%%%%%%%%%%%%%%%%%%%%%%%%%%%%%%%%%%%%%%%%%%%%%%%%%%%%%%%%%%%%%%%%%%%%%%%
\section*{Discussion}
The theorem is proved by a direct application of the smoothness
inequality together with the unbiasedness and variance assumptions.
The periodic cosine schedule does not affect the algebraic steps as long
as the per‑iteration stepsize stays within the interval
$[\eta_{\min},\eta_{\max}]$.  The only place where the schedule enters the
analysis is the bounding of $\sum_{t=0}^{T-1}\eta_t$ from below by
$\eta_{\min}T$ and $\sum_{t=0}^{T-1}\eta_t^{2}$ from above by
$\eta_{\max}\sum_{t=0}^{T-1}\eta_t\le\eta_{\max}^{2}T$.
Choosing $M_k$ to grow geometrically (the original SGDR proposal)
makes $\eta_{\max}$ effectively decay like $1/\sqrt{T}$, yielding the
optimal $\mathcal{O}(1/\sqrt{T})$ rate for non‑convex stochastic
optimization.

%%%%%%%%%%%%%%%%%%%%%%%%%%%%%%%%%%%%%%%%%%%%%%%%%%%%%%%%%%%%%%%%%%%%%%%%%%%%%%%
\section*{Preliminaries (Definitions and Lemmas)}
\begin{lemma}[Descent Lemma]\label{lem:descent}
If $f$ is $L$‑smooth, then for any $x$ and any vector $d$,
\[
f(x-d)\le f(x)-\langle\nabla f(x),d\rangle+\frac{L}{2}\|d\|^{2}.
\tag{8}
\]
\end{lemma}
\begin{proof}
Immediate from definition \eqref{3} with $y=x-d$.
\end{proof}

\begin{lemma}[Variance Decomposition]\label{lem:var}
For any random vector $g$,
\[
\mathbb{E}\|g\|^{2}= \|\mathbb{E}g\|^{2}
                     +\mathbb{E}\|g-\mathbb{E}g\|^{2}.
\tag{9}
\]
\end{lemma}
\begin{proof}
Expand $\|g\|^{2}=\|g-\mathbb{E}g+\mathbb{E}g\|^{2}$ and take expectation.
\end{proof}

%%%%%%%%%%%%%%%%%%%%%%%%%%%%%%%%%%%%%%%%%%%%%%%%%%%%%%%%%%%%%%%%%%%%%%%%%%%%%%%
\section*{Main Proof (Step‑by‑Step Derivation)}
We prove Theorem~\ref{thm:sgdr}.  All expectations are taken w.r.t. the
randomness up to time $t$ unless otherwise noted.

\noindent\textbf{Step 1.} Apply Lemma~\ref{lem:descent} with
$d=\eta_t g_t$ and $x=w_t$:
\begin{align}
f(w_{t+1})
&\le f(w_t)-\eta_t\langle\nabla f(w_t),g_t\rangle
     +\frac{L}{2}\eta_t^{2}\|g_t\|^{2}.
\tag{10}
\end{align}
\textbf{[Step 1]}

\noindent\textbf{Step 2.} Take conditional expectation w.r.t. $w_t$.
Using unbiasedness \eqref{4}:
\[
\mathbb{E}\bigl[\langle\nabla f(w_t),g_t\rangle\mid w_t\bigr]
   =\|\nabla f(w_t)\|^{2}.
\tag{11}
\]
\textbf{[Step 2]}

\noindent\textbf{Step 3.} Apply Lemma~\ref{lem:var} together with the
variance bound \eqref{5}:
\begin{align}
\mathbb{E}\bigl[\|g_t\|^{2}\mid w_t\bigr]
&= \|\nabla f(w_t)\|^{2}
   +\mathbb{E}\bigl[\|g_t-\nabla f(w_t)\|^{2}\mid w_t\bigr]
 \le \|\nabla f(w_t)\|^{2}+\sigma^{2}.
\tag{12}
\end{align}
\textbf{[Step 3]}

\noindent\textbf{Step 4.} Insert \eqref{11} and \eqref{12} into the
conditional expectation of \eqref{10}:
\begin{align}
\mathbb{E}\!\bigl[f(w_{t+1})\mid w_t\bigr]
&\le f(w_t)-\eta_t\|\nabla f(w_t)\|^{2}
   +\frac{L}{2}\eta_t^{2}\bigl(\|\nabla f(w_t)\|^{2}+\sigma^{2}\bigr).
\tag{13}
\end{align}
\textbf{[Step 4]}

\noindent\textbf{Step 5.} Rearrange \eqref{13}:
\begin{align}
\eta_t\Bigl(1-\frac{L}{2}\eta_t\Bigr)\|\nabla f(w_t)\|^{2}
&\le f(w_t)-\mathbb{E}\!\bigl[f(w_{t+1})\mid w_t\bigr]
   +\frac{L\sigma^{2}}{2}\eta_t^{2}.
\tag{14}
\end{align}
\textbf{[Step 5]}

\noindent\textbf{Step 6.} By assumption \eqref{6}, $\eta_t\le\eta_{\max}\le1/L$,
hence $1-\frac{L}{2}\eta_t\ge\frac12$.  Therefore
\[
\frac12\,\eta_t\|\nabla f(w_t)\|^{2}
\le f(w_t)-\mathbb{E}\!\bigl[f(w_{t+1})\mid w_t\bigr]
   +\frac{L\sigma^{2}}{2}\eta_t^{2}.
\tag{15}
\]
\textbf{[Step 6]}

\noindent\textbf{Step 7.} Take total expectation and sum \eqref{15} from
$t=0$ to $T-1$:
\begin{align}
\frac12\sum_{t=0}^{T-1}\eta_t\,
      \mathbb{E}\!\bigl[\|\nabla f(w_t)\|^{2}\bigr]
&\le \mathbb{E}[f(w_0)]-\mathbb{E}[f(w_T)]
   +\frac{L\sigma^{2}}{2}\sum_{t=0}^{T-1}\eta_t^{2}.
\tag{16}
\end{align}
\textbf{[Step 7]}

\noindent\textbf{Step 8.} Drop the non‑negative term $\mathbb{E}[f(w_T)]$ and
use $\Delta_f = f(w_0)-\inf f$:
\[
\frac12\sum_{t=0}^{T-1}\eta_t\,
      \mathbb{E}\!\bigl[\|\nabla f(w_t)\|^{2}\bigr]
\le \Delta_f
   +\frac{L\sigma^{2}}{2}\sum_{t=0}^{T-1}\eta_t^{2}.
\tag{17}
\]
\textbf{[Step 8]}

\noindent\textbf{Step 9.} Lower–bound the left‑hand side by using the
minimal stepsize:
\[
\sum_{t=0}^{T-1}\eta_t\,
      \mathbb{E}\!\bigl[\|\nabla f(w_t)\|^{2}\bigr]
\ge \eta_{\min}\sum_{t=0}^{T-1}
      \mathbb{E}\!\bigl[\|\nabla f(w_t)\|^{2}\bigr].
\tag{18}
\]
\textbf{[Step 9]}

\noindent\textbf{Step 10.} Upper–bound the squared stepsizes using the
maximal stepsize:
\[
\sum_{t=0}^{T-1}\eta_t^{2}
\le \eta_{\max}\sum_{t=0}^{T-1}\eta_t
\le \eta_{\max}^{2}T .
\tag{19}
\]
\textbf{[Step 10]}

\noindent\textbf{Step 11.} Combine \eqref{17}, \eqref{18} and \eqref{19}:
\[
\eta_{\min}\sum_{t=0}^{T-1}
      \mathbb{E}\!\bigl[\|\nabla f(w_t)\|^{2}\bigr]
\le 2\Delta_f + L\sigma^{2}\eta_{\max}^{2}T .
\tag{20}
\]
\textbf{[Step 11]}

\noindent\textbf{Step 12.} Divide both sides of \eqref{20} by $\eta_{\min}T$:
\[
\frac{1}{T}\sum_{t=0}^{T-1}
      \mathbb{E}\!\bigl[\|\nabla f(w_t)\|^{2}\bigr]
\le \frac{2\Delta_f}{\eta_{\min}T}
   +\frac{L\sigma^{2}\eta_{\max}^{2}}{\eta_{\min}}}{\displaystyle
      \frac{L\sigma^{2}\eta_{\max}^{2}}{\eta_{\min}}}
   =\frac{2\Delta_f}{\eta_{\min}T}
    +\frac{L\sigma^{2}\eta_{\max}}{\eta_{\min}} .
\tag{21}
\]
\textbf{[Step 12]}

Equation \eqref{21} is exactly the bound \eqref{7} claimed in
Theorem~\ref{thm:sgdr}.  $\square$

%%%%%%%%%%%%%%%%%%%%%%%%%%%%%%%%%%%%%%%%%%%%%%%%%%%%%%%%%%%%%%%%%%%%%%%%%%%%%%%
\section*{Rate Derivation (With Constants)}
If the period lengths are chosen as
\[
M_k = M_0\,T_{\rm mult}^{\,k},\qquad k=0,1,\dots,
\]
the total number of iterations after $K$ restarts is
\[
T = \sum_{k=0}^{K} M_k
  = M_0\frac{T_{\rm mult}^{\,K+1}-1}{T_{\rm mult}-1}.
\]
The maximal stepsize inside the $k$‑th period equals
\[
\eta_{\max}^{(k)} = \eta_{\max}^{(0)}\;T_{\rm mult}^{\,-k/2},
\]
i.e. it decays like $1/\sqrt{M_k}$.  Substituting
$\eta_{\max}=c/\sqrt{T}$ into \eqref{7} yields
\[
\frac{1}{T}\sum_{t=0}^{T-1}\mathbb{E}\|\nabla f(w_t)\|^{2}
\;\le\;
\frac{2\Delta_f}{\eta_{\min}T}
   +\frac{Lc\sigma^{2}}{\eta_{\min}\sqrt{T}}
   =\mathcal{O}\!\bigl(T^{-1/2}\bigr).
\]
Thus SGDR attains the optimal stochastic non‑convex rate.

%%%%%%%%%%%%%%%%%%%%%%%%%%%%%%%%%%%%%%%%%%%%%%%%%%%%%%%%%%%%%%%%%%%%%%%%%%%%%%%
\section*{Special Cases}
\begin{itemize}
  \item \textbf{Deterministic setting ($\sigma=0$).}  Equation \eqref{7}
    reduces to $\frac{1}{T}\sum\|\nabla f(w_t)\|^{2}\le
    \frac{2\Delta_f}{\eta_{\min}T}$, i.e. a $O(1/T)$ convergence of the
    gradient norm, matching the standard result for deterministic
    gradient descent with diminishing stepsizes.
  \item \textbf{Constant stepsize ($\eta_{\min}=\eta_{\max}=\eta$).}
    The bound becomes
    $\frac{1}{T}\sum\mathbb{E}\|\nabla f(w_t)\|^{2}
      \le \frac{2\Delta_f}{\eta T}+L\sigma^{2}\eta$,
    which is the classic SGD guarantee.
  \item \textbf{No restart ($M_k\equiv\infty$).}
    The schedule collapses to the ordinary cosine annealing without
    reset; the proof still holds because the only required property is
    $\eta_t\in[\eta_{\min},\eta_{\max}]$.
\end{itemize}

%%%%%%%%%%%%%%%%%%%%%%%%%%%%%%%%%%%%%%%%%%%%%%%%%%%%%%%%%%%%%%%%%%%%%%%%%%%%%%%
\end{document}